\section{Marks}

Marks are graphical elements placed at the intersections of the board to highlight specific positions, indicate points of interest, or number sequences without the need to add stones. \pgfgoname\ offers versatile commands and options to customize these marks, ranging from predefined shapes to custom labels.

\begin{command}{\mark\opt{\oarg{options}}\marg{intersection}}
	
	Draws an individual mark at the specified intersection (e.g., "C3"). Options allow defining its shape, color, and other attributes.
\end{command}

\begin{command}{\marks\opt{\oarg{options}}\marg{intersection \opt{, intersection, ...}}}
	
	Similar to |\mark|, but applies the same mark to multiple intersections using the coordinate syntax.
\end{command}

\begin{key}{/pgfgo/mark = \meta{options} (default cross)}
	
	Defines the type of mark to be placed, with a cross ("cross") as the default value if no other option is specified.
	
	\begin{codeexample}[]
		\pgfgoset{goban = {5, partial = {from = B2 to D4}}}
		
		\begin{goban}
			\stones[b]{C2,D3}
			\stone[w]{B4}
		\end{goban}
		
		\begin{goban}
			\stones[b]{C2,D3}
			\stone[w]{B4}
			\mark{C3}
		\end{goban}
	\end{codeexample}
\end{key}

\begin{key}{/pgfgo/mark/circle}
	
	Places a hollow circle as a mark.
	
	\begin{codeexample}[]
		\pgfgoset{goban = {5, partial = {from = B2 to D4}}}
		
		\begin{goban}
			\stones[b]{C2,D3}
			\stone[w]{B4}
		\end{goban}
		
		\begin{goban}
			\stones[b]{C2,D3}
			\stone[w]{B4}
			\mark[circle]{C3}
		\end{goban}
	\end{codeexample}
\end{key}

\begin{key}{/pgfgo/mark/triangle}
	
	Draws a hollow triangle.
	
	\begin{codeexample}[]
		\pgfgoset{goban = {5, partial = {from = B2 to D4}}}
		
		\begin{goban}
			\stones[b]{C2,D3}
			\stone[w]{B4}
		\end{goban}
		
		\begin{goban}
			\stones[b]{C2,D3}
			\stone[w]{B4}
			\mark[triangle]{C3}
		\end{goban}
	\end{codeexample}
\end{key}

\begin{key}{/pgfgo/mark/square}
	
	Places a hollow square.
	
	\begin{codeexample}[]
		\pgfgoset{goban = {5, partial = {from = B2 to D4}}}
		
		\begin{goban}
			\stones[b]{C2,D3}
			\stone[w]{B4}
		\end{goban}
		
		\begin{goban}
			\stones[b]{C2,D3}
			\stone[w]{B4}
			\mark[square]{C3}
		\end{goban}
	\end{codeexample}
\end{key}

\begin{key}{/pgfgo/mark/cross}
	
	Draws a cross.
	
	\begin{codeexample}[]
		\pgfgoset{goban = {5, partial = {from = B2 to D4}}}
		
		\begin{goban}
			\stones[b]{C2,D3}
			\stone[w]{B4}
		\end{goban}
		
		\begin{goban}
			\stones[b]{C2,D3}
			\stone[w]{B4}
			\mark[cross]{C3}
		\end{goban}
	\end{codeexample}
\end{key}

\begin{key}{/pgfgo/mark/filled square}
	
	Places a filled square.
	
	\begin{codeexample}[]
		\pgfgoset{goban = {5, partial = {from = B2 to D4}}}
		
		\begin{goban}
			\stones[b]{C2,D3}
			\stone[w]{B4}
		\end{goban}
		
		\begin{goban}
			\stones[b]{C2,D3}
			\stone[w]{B4}
			\mark[filled square]{C3}
		\end{goban}
	\end{codeexample}
\end{key}

\begin{key}{/pgfgo/mark/filled circle}
	
	Draws a filled circle.
	
	\begin{codeexample}[]
		\pgfgoset{goban = {5, partial = {from = B2 to D4}}}
		
		\begin{goban}
			\stones[b]{C2,D3}
			\stone[w]{B4}
		\end{goban}
		
		\begin{goban}
			\stones[b]{C2,D3}
			\stone[w]{B4}
			\mark[filled circle]{C3}
		\end{goban}
	\end{codeexample}
\end{key}

\begin{key}{/pgfgo/mark/t-circle}
	
	Places a small circle with a thicker border, specifically designed for marking territories. It can also be written |tcircle|.
	
	\begin{codeexample}[]
		\pgfgoset{goban = {5, partial = {from = B2 to D4}}}
		
		\begin{goban}
			\stones[b]{C2,D3}
			\stone[w]{B4}
		\end{goban}
		
		\begin{goban}
			\stones[b]{C2,D3}
			\stone[w]{B4}
			\mark[t-circle]{C3}
		\end{goban}
	\end{codeexample}
\end{key}

\begin{key}{/pgfgo/mark/t-square}
	
	Draws a small square with a thick border, similar to "t-circle," used to mark territories. It can also be written |tsquare|.
	
	\begin{codeexample}[]
		\pgfgoset{goban = {5, partial = {from = B2 to D4}}}
		
		\begin{goban}
			\stones[b]{C2,D3}
			\stone[w]{B4}
		\end{goban}
		
		\begin{goban}
			\stones[b]{C2,D3}
			\stone[w]{B4}
			\mark[t-square]{C3}
		\end{goban}
	\end{codeexample}
\end{key}

\begin{key}{/pgfgo/mark/hatched lines}
	
	Places a pattern of crossed (hatched) lines.
	
	\begin{codeexample}[]
		\pgfgoset{goban = {5, partial = {from = B2 to D4}}}
		
		\begin{goban}
			\stones[b]{C2,D3}
			\stone[w]{B4}
		\end{goban}
		
		\begin{goban}
			\stones[b]{C2,D3}
			\stone[w]{B4}
			\mark[hatched lines]{C3}
		\end{goban}
	\end{codeexample}
\end{key}

\begin{key}{/pgfgo/mark/sequence=\meta{options} (default from = 1)}
	Allows numbering a series of marks in order.
\end{key}

\begin{key}{/pgfgo/mark/sequence/use}
	Displays the |sequence| mark on the board. It is the same as |sequence = true|.
\end{key}

\begin{key}{/pgfgo/mark/sequence/false}
	Turns the sequence off (for example, one turned on with |\pgfgoset|).
\end{key}

\begin{key}{/pgfgo/mark/sequence/format = \meta{format}}
	Turns the sequence on and chooses the number format. It is the same as |sequence, label = {format = ...}| (see |label/format|).
\end{key}

\begin{key}{/pgfgo/mark/sequence/from=\meta{number}}
	
	Sets the starting number of a numbered mark sequence, allowing continuation from a specific point instead of 1.
	
	\begin{codeexample}[]
		\pgfgoset{goban = {5, partial = {from = B2 to D4}}}
		
		\begin{goban}
			\stones[b]{C2,D3}
			\stone[w]{B4}
		\end{goban}
		
		\begin{goban}
			\stones[b]{C2,D3}
			\stone[w]{B4}
			\marks[sequence = {use, from = 3}]{C3,B3,C4}
		\end{goban}
	\end{codeexample}
\end{key}

\begin{key}{/pgfgo/mark/custom=\meta{token list}}
	
	Allows defining a custom mark using text or symbols (e.g., "!" or "?") and even images.
	
	\begin{codeexample}[]
		\pgfgoset{goban = {5, partial = {from = B2 to D4}}}
		
		\begin{goban}
			\stones[b]{C2,D3}
			\stone[w]{B4}
		\end{goban}
		
		\begin{goban}
			\stones[b]{C2,D3}
			\stone[w]{B4}
			\mark[custom = {!}]{C3}
		\end{goban}
		
		\begin{goban}
			\stones[b]{C2,D3}
			\stone[w]{B4}
			\mark[custom = {\tikz{\node[rounded corners, draw, inner sep = 1em, red]{\dots};}}]{C3}
		\end{goban}
	\end{codeexample}
\end{key}

\begin{key}{/pgfgo/mark/line = \meta{options}}
	Controls the outline of the mark.
\end{key}

\begin{key}{/pgfgo/mark/line/color = \meta{color}}
	Defines the outline color of the mark.
	
	\begin{codeexample}[]
		\pgfgoset{goban = {5, partial = {from = B2 to D4}}}
		
		\begin{goban}
			\stones[b]{C2,D3}
			\stone[w]{B4}
		\end{goban}
		
		\begin{goban}
			\stones[b]{C2,D3}
			\stone[w]{B4}
			\mark[triangle, line = {color = red}]{C3}
		\end{goban}
	\end{codeexample}
\end{key}

\begin{key}{/pgfgo/mark/line/width = \meta{dimension}}
	
	Adjusts the thickness of the mark’s outline.
	
	\begin{codeexample}[]
		\pgfgoset{goban = {5, partial = {from = B2 to D4}}}
		
		\begin{goban}
			\stones[b]{C2,D3}
			\stone[w]{B4}
			\mark[triangle]{C3}
		\end{goban}
		
		\begin{goban}
			\stones[b]{C2,D3}
			\stone[w]{B4}
			\mark[triangle, line = {width = 0.1pt}]{C3}
		\end{goban}
	\end{codeexample}
\end{key}

\begin{key}{/pgfgo/mark/fill = \meta{options}}
	Controls the fill of the filled marks (|filled square|, |filled circle| and |hatched lines|).
\end{key}

\begin{key}{/pgfgo/mark/fill/true}
	Filled marks use |fill/color|. Without |fill = true|, they are filled in gray.
\end{key}

\begin{key}{/pgfgo/mark/fill/false}
	Filled marks are filled in gray again. It is the initial value.
\end{key}

\begin{key}{/pgfgo/mark/fill/color = \meta{color}}
	
	Sets the fill color of a mark. It is only used with |fill = true|.
	
	\begin{codeexample}[]
		\pgfgoset{goban = {5, partial = {from = B2 to D4}}}
		
		\begin{goban}
			\stones[b]{C2,D3}
			\stone[w]{B4}
		\end{goban}
		
		\begin{goban}
			\stones[b]{C2,D3}
			\stone[w]{B4}
			\mark[filled circle, fill = {true, color = red}]{C3}
		\end{goban}
	\end{codeexample}
\end{key}

\begin{key}{/pgfgo/mark/label = \meta{option}}
	Controls the appearance of labels for marks that have them.
\end{key}

\begin{key}{/pgfgo/mark/label/color = \meta{color}}
	
	Controls the color of the text or symbol in custom marks (e.g., with |custom|).
	
	\begin{codeexample}[]
		\pgfgoset{goban = {5, partial = {from = B2 to D4}}}
		
		\begin{goban}
			\stones[b]{C2,D3}
			\stone[w]{B4}
		\end{goban}
		
		\begin{goban}
			\stones[b]{C2,D3}
			\stone[w]{B4}
			\mark[custom = {!}, label = {color = red}]{C3}
		\end{goban}
	\end{codeexample}
\end{key}

\begin{key}{/pgfgo/mark/label/font = \meta{font commands}}
	Font of the labels of marks (|sequence| and |custom|).

	\begin{codeexample}[]
		\pgfgoset{goban = {5, partial = {from = B2 to D4}}}

		\begin{goban}
			\stones[b]{C2,D3}
			\stone[w]{B4}
			\marks[sequence, label = {font = \itshape}]{C3,B3,C4}
		\end{goban}
	\end{codeexample}
\end{key}

\begin{key}{/pgfgo/mark/label/format=\meta{options} (initially arabic)}
	
	Controls the numbering style of enumeration labels (e.g., |sequence|), with options such as Arabic numerals, letters, or Roman numerals.
\end{key}

\begin{key}{/pgfgo/mark/label/format/arabic}
	
	Uses Arabic numerals for sequences.
	
	\begin{codeexample}[]
		\pgfgoset{goban = {5, partial = {from = B2 to D4}}}
		
		\begin{goban}
			\stones[b]{C2,D3}
			\stone[w]{B4}
		\end{goban}
		
		\begin{goban}
			\stones[b]{C2,D3}
			\stone[w]{B4}
			\marks[label = {format = arabic}, sequence]{C3,B3,C4}
		\end{goban}
	\end{codeexample}
\end{key}

\begin{key}{/pgfgo/mark/label/format/alph}
	
	Uses lowercase letters for the sequence.
	
	\begin{codeexample}[]
		\pgfgoset{goban = {5, partial = {from = B2 to D4}}}
		
		\begin{goban}
			\stones[b]{C2,D3}
			\stone[w]{B4}
		\end{goban}
		
		\begin{goban}
			\stones[b]{C2,D3}
			\stone[w]{B4}
			\marks[label = {format = alph}, sequence]{C3,B3,C4}
		\end{goban}
	\end{codeexample}
\end{key}

\begin{key}{/pgfgo/mark/label/format/Alph}
	
	Uses uppercase letters.
	
	\begin{codeexample}[]
		\pgfgoset{goban = {5, partial = {from = B2 to D4}}}
		
		\begin{goban}
			\stones[b]{C2,D3}
			\stone[w]{B4}
		\end{goban}
		
		\begin{goban}
			\stones[b]{C2,D3}
			\stone[w]{B4}
			\marks[label = {format = Alph}, sequence]{C3,B3,C4}
		\end{goban}
	\end{codeexample}
\end{key}

\begin{key}{/pgfgo/mark/label/format/roman}
	
	Uses lowercase Roman numerals.
	
	\begin{codeexample}[]
		\pgfgoset{goban = {5, partial = {from = B2 to D4}}}
		
		\begin{goban}
			\stones[b]{C2,D3}
			\stone[w]{B4}
		\end{goban}
		
		\begin{goban}
			\stones[b]{C2,D3}
			\stone[w]{B4}
			\marks[label = {format = roman}, sequence]{C3,B3,C4}
		\end{goban}
	\end{codeexample}
\end{key}

\begin{key}{/pgfgo/mark/label/format/Roman}
	
	Uses uppercase Roman numerals.
	
	\begin{codeexample}[]
		\pgfgoset{goban = {5, partial = {from = B2 to D4}}}
		
		\begin{goban}
			\stones[b]{C2,D3}
			\stone[w]{B4}
		\end{goban}
		
		\begin{goban}
			\stones[b]{C2,D3}
			\stone[w]{B4}
			\marks[label = {format = Roman}, sequence]{C3,B3,C4}
		\end{goban}
	\end{codeexample}
\end{key}

\begin{key}{/pgfgo/mark/black}
	
	Applies the properties of the first player (black) to the mark.
	
	\begin{codeexample}[]
		\pgfgoset{goban = {5, partial = {from = B2 to D4}}}
		
		\begin{goban}
			\stones[b]{C2,D3}
			\stone[w]{B4}
			\mark[triangle]{C3}
		\end{goban}
		
		\begin{goban}
			\stones[b]{C2,D3}
			\stone[w]{B4}
			\mark[black, t-square]{C3}
		\end{goban}
	\end{codeexample}
\end{key}

\begin{key}{/pgfgo/mark/white}
	
	Applies the properties of the second player (white) to the mark.
\end{key}

\begin{key}{/pgfgo/mark/neutral}
	
	Keeps the mark unassociated with any player.
\end{key}

\begin{key}{/pgfgo/mark/b}
	Short for |black|.
\end{key}

\begin{key}{/pgfgo/mark/w}
	Short for |white|.
\end{key}

\begin{key}{/pgfgo/mark/n}
	Short for |neutral|.
\end{key}

\begin{key}{/pgfgo/mark/forget}
	This mark is not stored in the board’s memory, even if the |goban| uses |remember|.
\end{key}

\begin{key}{/pgfgo/mark/remember}
	This mark is stored in the board’s memory, even if the |goban| uses |forget|.
\end{key}